\Askhsh{34733}{4}
\begin{ylh}{1}
    \item 2.2 Παρουσίαση Στατιστικών Δεδομένων
\end{ylh}
Σε έναν κινηματογράφο τα εισιτήρια που έκοψαν τέσσερεις ταινίες τον μήνα Νοέμβριο φαίνονται στον επόμενο πίνακα:
\begin{center}
    \begin{myhtblr}{}
        
        Ταινία 1 & 150 \\
        
        Ταινία 2 & 130 \\
        
        Ταινία 3 & 70  \\
        
        Ταινία 4 & 50  \\
        
    \end{myhtblr}
\end{center}

\begin{alist}
    \item Ένας μαθητής κατασκεύασε το επόμενο πολύγωνο σχετικών συχνοτήτων \% για τους θεατές που παρακολούθησαν τον μήνα Νοέμβριο τις παραπάνω ταινίες:
    \begin{center}
        \begin{tikzpicture}
        \begin{axis}[
            width=12cm, height=6cm,
            axis lines=box,
            xtick={1,2,3,4},
            xticklabels={Ταινία 1, Ταινία 2, Ταινία 3, Ταινία 4},
            ymin=0, ymax=170,
            ytick=\empty,
            ymajorgrids=true,
            grid style={solid},
            xtick style={draw=none},
        ]
            \addplot[mark=*, mark options={fill=maincolor, draw=maincolor}, thick, maincolor] coordinates {(1,150) (2,130) (3,70) (4,50)};
        \end{axis}
        \end{tikzpicture}
    \end{center}
    Να αιτιολογήσετε αν το πολύγωνο αυτό είναι σωστό ή λανθασμένο.
    \monades{8}

    \item Να συμπληρώσετε τον παρακάτω πίνακα:
    \begin{center}
\begin{tblr}{Q[f] Q[f]}
        \begin{mytblr}{}
            \textbf{Ταινία} & \textbf{Συχνότητα} & \textbf{Σχετική συχνότητα $f_i$} \\
            Ταινία 1        & 150                &                                  \\ 
            Ταινία 2        & 130                &                                  \\
            Ταινία 3        & 70                 &                                  \\
            Ταινία 4        & 50                 &                                  \\
            \textbf{Σύνολο} &                    &                                  \\
        \end{mytblr}&
\begin{tikzpicture}[scale=1.2]
            \def\radius{2.5cm}
            % Diameters
            \draw[maincolor] (-\radius,0) -- (\radius,0);
            \draw[secondarycolor] (0,-\radius) -- (0,\radius);
            % Center point
            \node[circle, fill=maincolor, inner sep=1pt] at (0,0) {};
            % Circle and dots
            \draw (0,0) circle (\radius);
            \foreach \i in {0,5,...,355}
            {
                \node[circle, fill=black, inner sep=0.8pt] at (\i:\radius) {};
            }
        \end{tikzpicture}
    \end{tblr}
\end{center}
    \monades{9}

    \item Να υπολογίσετε σε μοίρες τη γωνία που θα αντιστοιχεί σε ένα κυκλικό διάγραμμα σχετικών συχνοτήτων για καθεμία από τις παραπάνω ταινίες. Στη συνέχεια να κατασκευάσετε ένα κυκλικό διάγραμμα σχετικών συχνοτήτων, χρησιμοποιώντας τον παραπάνω κυκλικό δίσκο, ο οποίος είναι χωρισμένος σε τόξα $5^\circ$.
    \monades{8}
\end{alist}

